\subsection{Светодиоды и диагностическая индикация}
 
После подключения к дрону полезно получить первый видимый результат без взлета. Для этого удобно использовать светодиодную индикацию: она безопасна, мгновенно показывает состояние программы и помогает отлаживать реакции на данные датчиков.
 
\subsubsection*{Зачем начинать со светодиодов?}
\begin{itemize}
    \item Светодиоды дают быстрый и наглядный результат без полета.
    \item Через цвет удобно кодировать состояния: норма, предупреждение, ошибка.
    \item Эта тема естественно связывает базовый Python с реальным API квадрокоптера.
\end{itemize}
 
\subsubsection*{Базовая команда}
Для управления индикацией используется метод \texttt{led\_control}. Чаще всего в учебных задачах применяется идентификатор \texttt{255}, который означает управление всей подсветкой сразу.
 
\begin{codefile}[python]{python/examples/03_flight_control/LED_hight_change.py}
\end{codefile}
 
\begin{infobox}[Как читать вызов]
\texttt{drone.led\_control(led\_id=255, r, g, b)} устанавливает цвет подсветки. Компоненты \texttt{r}, \texttt{g}, \texttt{b} задают интенсивность красного, зеленого и синего каналов.
\end{infobox}
 
\subsubsection*{Минимальная палитра}
\begin{itemize}
    \item \texttt{(1, 0, 0)} --- красный: ошибка, авария, низкий заряд.
    \item \texttt{(0, 1, 0)} --- зеленый: норма, готовность.
    \item \texttt{(0, 0, 1)} --- синий: ожидание, служебный режим.
    \item \texttt{(1, 1, 0)} --- желтый: предупреждение.
    \item \texttt{(0, 0, 0)} --- выключено.
\end{itemize}
 
\begin{warnbox}[Безопасный режим]
На этом этапе не нужно взлетать. Все упражнения со светодиодами можно выполнять на столе, проверяя только подключение, чтение телеметрии и корректную реакцию программы.
\end{warnbox}
 
\subsubsection*{Светодиоды как первый дебаг-интерфейс}
Даже если консоль недоступна или оператор смотрит не в терминал, цветовая индикация позволяет сразу понять состояние программы:
\begin{itemize}
    \item зеленый --- соединение установлено;
    \item желтый --- состояние требует внимания;
    \item красный --- нужно остановить сценарий и проверить проблему.
\end{itemize}
 
\begin{taskbox}[Первая индикация]
Напишите скрипт, который:
\begin{enumerate}
    \item подключается к дрону;
    \item включает зеленую подсветку на 2 секунды;
    \item затем выключает ее;
    \item корректно закрывает соединение.
\end{enumerate}
\end{taskbox}
 
\begin{taskbox}[Светофор батареи]
Используя \texttt{get\_battery\_status()}, реализуйте простую логику:
\begin{itemize}
    \item выше 4.0 В --- зеленый;
    \item от 3.7 до 4.0 В --- желтый;
    \item ниже 3.7 В --- красный.
\end{itemize}
\end{taskbox}
 
\begin{taskbox}[Индикация по высоте]
Измените пример \texttt{LED\_hight\_change.py}, чтобы цвет зависел от высоты по вашим порогам. Например: низко --- синий, рабочая высота --- зеленый, слишком высоко --- красный.
\end{taskbox}
 
\subsubsection*{Что это готовит дальше}
После этого раздела можно переходить к полету значительно увереннее: ученик уже умеет подключаться к дрону, вызывать методы API и связывать данные датчиков с понятной реакцией устройства.
